\section{Actual interior test histories and the recursive monomial leaf}

An original crossing lower diagram has the same actual initial lower
record but a different entire tail, containing an internal identity empty
blowup.  Its admissibility must therefore be supplied from the original
data, not assumed as a lower-history comparison.  We first establish that
real test-history preservation.  The actual monomial branch of the full
recursive invariant is then transported along these original histories.
This is additional explicit indexed-history bookkeeping for Wlodarczyk's
Step One and Bierstone--Milman's trivial blowup convention.

\begin{lemma}[Normal crossings ignore an inserted empty boundary entry]
\label{hir:original-empty-boundary-centre-normal-crossings}
\lean{Hironaka.HironakaOriginalCrossingValueTransport.hasNormalCrossings_insertIdx_empty_iff}
\leanok
\source{wlodarczyk05-hironaka,bierstone-milman08-functoriality}
\uses{hir:snc-basic}
For every actual ordered boundary list $E$, actual subset $C$ and index
$j\leq |E|$, $C$ has normal crossings with
$E.\operatorname{insertIdx}(j,\varnothing)$ if and only if it has normal
crossings with $E$.  No boundary entry is deleted from either history.
\end{lemma}
\begin{proof}
At a point $x\in C$, retain the actual adapted chart and coordinate
support of $C$.  A member of the inserted list through $x$ cannot be its
empty entry, so it is an actual old member.  Choose its original list
index and use that member's actual coordinate support; for a member
missing $x$ choose empty support.  Each required coordinate formula and
one-coordinate cardinality is precisely the retained old member's
formula.  The inserted entry has no required formula or cardinality,
because no point belongs to it.  This constructs adapted coordinates for
the entire inserted indexed list.  Conversely every old member occurs
in the actual inserted list.  At $x$, choose its actual index there and
retain that coordinate support, with the same chart and $C$ support.
Its formula and cardinality are the supplied inserted-list ones.  These
constructions give both directions pointwise, and hence for all $x\in C$.
They use genuine membership equality of the nonempty members; no global
index-preimage injectivity or new normal-crossings certificate is needed.
\end{proof}

\begin{lemma}[Actual test towers under an initial empty boundary insertion]
\label{hir:original-empty-boundary-test-history}
\lean{Hironaka.HironakaOriginalCrossingValueTransport.isAdmissibleCentre_insertEmpty_iff,Hironaka.HironakaOriginalCrossingValueTransport.isGeoTestTower_insertEmpty_iff}
\leanok
\source{wlodarczyk05-hironaka,bierstone-milman08-functoriality}
\uses{hir:original-empty-boundary-centre-normal-crossings,hir:original-invariant-empty-actual}
For any marked ideal $J$ and $j\leq |J.E|$, any actual tower $R$ is a
test tower of $\operatorname{insertEmpty}_j(J)$ if and only if it is a
test tower of $J$.  The actual admissibility of every centre also agrees.
\end{lemma}
\begin{proof}
Inserting the empty member changes neither the carrier, coefficient
ideal, mark nor cosupport.  A centre has the same closed submanifold and
cosupport conditions, and its normal-crossings condition agrees by the
previous lemma.  Thus its entire admissibility condition agrees.
Induct on the actual last steps of $R$.  The empty tower is a test tower
on both sides.  At a last step, actual transformation commutes with the
same indexed insertion; the index remains valid because the whole
boundary has gained exactly one entry at each previous step.  The
induction hypothesis gives testness of the predecessor and the proved
admissibility equivalence gives its actual last centre.  This proves
both directions for the entire tower, with all original blowup data kept.
\end{proof}

\begin{theorem}[Genuine original testness across the internal identity empty step]
\label{hir:original-interior-empty-test-history}
\lean{Hironaka.HironakaOriginalCrossingValueTransport.isGeoTestTower_interior_empty_iff}
\leanok
\source{wlodarczyk05-hironaka,bierstone-milman08-functoriality}
\uses{hir:original-empty-boundary-test-history,hir:original-interior-empty-geometric-dictionary,hir:restricted-equivalence,hir:wellformed-tower}
Let $M$ be originally well formed and $H$ an actual test tower of $M$.
For every actual tail $S$, $H\epsilon S$ is an actual test tower of $M$
if and only if $HS$ is, where $\epsilon$ is the genuine identity empty
blowup.  This includes an empty tail.
\end{theorem}
\begin{proof}
The actual transformed original record $K=M_H$ is well formed.  Its
genuine identity empty step is admissible, so $H\epsilon$ is an actual
test head.  The actual transformed record after that step is
$\operatorname{insertEmpty}_{|M.E|+|H|}(K_\epsilon^-)$, where
$K_\epsilon^-$ is the actual empty-step base with only its newly created
empty member removed.  The actual identity empty-blowup calculus gives
restricted equivalence of $K_\epsilon^-$ and $K$, genuine coherence of
both ideals, closedness of their common carrier, and equality of their
original whole boundary.  Restricted-equivalent coherent records have
the same actual admissible test towers.  Removing the new empty boundary
entry for testing is also legitimate by the previous lemma, with its
index bounded by the real whole empty-step base boundary.  Thus $S$ is
a test tail of the actual new head record exactly when it is a test tail
of $K$.  Testness of a concatenation is testness of its actual head and
tail; the already proved test heads $H$ and $H\epsilon$ now give precisely
the asserted equivalence.  No testness of the inserted lower tail is a
caller conclusion certificate.
\end{proof}

\begin{theorem}[The actual full-family monomial leaf in the same original history]
\label{hir:original-interior-empty-full-family-monomial-leaf}
\lean{Hironaka.HironakaOriginalCrossingValueTransport.nmOrd_interior_empty_nonempty,Hironaka.HironakaOriginalCrossingValueTransport.invFamily_interior_monomial_leaf,Hironaka.HironakaOriginalCrossingValueTransport.invFamily_interior_monomial_comparisons}
\leanok
\source{wlodarczyk05-hironaka,bierstone-milman08-functoriality}
\uses{hir:original-interior-empty-test-history,hir:original-interior-empty-geometric-dictionary,hir:original-invariant-empty-map-strict,hir:general-value-basic}
Let $n>0$, $k\geq0$, $M$ originally well formed, $H$ an actual test
history of $M$, $S$ a nonempty actual tail and $y$ an actual final point.
Assume that the actual nonmonomial order of $M_{HS}$ at $y$ is zero.
Then the actual nonmonomial order at $H\epsilon S,y$ is also zero and
\[
 \operatorname{invFamily}(n,k)(M,H\epsilon S,y)
 =\operatorname{shiftFamily}(k,|M.E|+|H|)
    \bigl(\operatorname{invFamily}(n,k)(M,HS,y)\bigr).
\]
For any two such actual zero-order points, strict order, weak order and
equality of these actual full invariant values are preserved and reflected.
The original initial ideal and entire history remain unre-based.
\end{theorem}
\begin{proof}
The nonempty-tail whole-record producer gives the actual transformed
record at $H\epsilon S$ as the old transformed record with the empty
entry inserted at the true position $j=|M.E|+|H|$, bounded by the old
whole boundary.  Actual nonmonomial order is unchanged by this insertion,
so both points lie in the same actual zero-order branch.  For $k=0$ the
actual family is constant and its shift is the identity.  For $k>0$ its
definition in this branch is its unchanged first coordinate zero and the
actual monomial invariant of its transformed original record.  That
monomial invariant is changed precisely by the proved indexed shift at
$j$.  The definition of shiftFamily on this branch is exactly that map,
giving the entire full-family equality.  The map is strictly increasing
on the full value type at every $k$, so applying it to two genuine
zero-order points preserves and reflects all three comparisons.  This
leaf computation uses no proposed arbitrary-invariant insertion axiom.
It does not replace the still-needed positive-order recursive crossing
transport or claim maximum existence or scheduling.
\end{proof}
